\section{Introduction}

When state-of-the-art language models achieve high scores on one trustworthiness benchmark, they systematically excel on others---even those measuring seemingly unrelated capabilities like mathematical reasoning and instruction following. Across the Open LLM Leaderboard's 4,561 models, benchmark scores correlate at $\rho = 0.80$--$0.87$, a pattern too strong and too consistent to be coincidental. This observation challenges a foundational assumption in LLM evaluation: that trustworthiness dimensions like truthfulness, robustness, and reliability are independent constructs requiring separate measurement.

The standard approach treats each trustworthiness dimension as a distinct target. Specialized benchmarks---TruthfulQA for factual accuracy, MMLU for knowledge, AdvGLUE for robustness---are developed, validated, and deployed independently. Practitioners evaluate models dimension-by-dimension, constructing scorecards that imply orthogonal capabilities. Yet the high cross-correlation suggests this framing may fundamentally mischaracterize how trustworthiness manifests in language models.

We identify a deeper problem: the correlation structure itself has been treated as a nuisance confound rather than an informative signal. Prior work attributes shared variance to model scale---larger models score higher on everything---and dismisses the residual as noise. But what if the residual is the signal? What if, after controlling for scale and training recency, a dominant latent factor persists that reflects something meaningful about model behavior?

This gap---the absence of a systematic factor-analytic framework for understanding cross-benchmark correlation in LLM trustworthiness---motivates our work. We draw on psychometric methodology, where similar ``positive manifold'' patterns led to the discovery of general intelligence ($g$-factor), to investigate whether an analogous latent construct underlies trustworthy LLM behavior.

Our key insight is that cross-benchmark correlation reflects a shared latent factor arising from \emph{representation stability}. Models with more stable internal representations produce consistent outputs across semantically equivalent inputs, simultaneously boosting performance on benchmarks measuring different surface-level capabilities. We term this latent factor \textbf{Generalized Representational Coherence (GRC)}.

Building on this insight, we make the following contributions:

\begin{enumerate}
    \item \textbf{Existence of residual factor:} Through large-scale meta-analysis ($N = 4{,}561$ models), we demonstrate that a dominant principal component (PC1, $\lambda_1 = 2.277$, explaining 60\% of residual variance) persists after controlling for $\log(\text{parameters})$ and release date, significantly exceeding permutation null thresholds ($p = 0.001$).

    \item \textbf{Mechanistic evidence:} We show that PC1 correlates positively with a Behavioral Stability Index (BSI), a proxy for representation consistency ($\rho = 0.405$, $p < 10^{-179}$), supporting the interpretation that GRC reflects stability rather than general capability.

    \item \textbf{Quasi-intervention evidence:} Instruction-tuning increases both BSI (Cohen's $d = 1.87$) and PC1 scores ($d = 1.99$) within matched base/instruct model pairs, suggesting that stability-enhancing training procedures causally improve GRC.

    \item \textbf{Prospective validity:} The frozen PC1 weights generalize to 5 of 6 holdout trustworthiness benchmarks (loadings $\geq 0.3$), demonstrating that GRC captures a genuine latent dimension rather than benchmark-specific artifacts.
\end{enumerate}

These findings suggest that evaluating trustworthiness as a single, latent factor---rather than a collection of independent dimensions---may better reflect the underlying structure of model capabilities.
